\begin{frame}
	\begin{center}
		\Huge Engine
	\end{center}
\end{frame}


\begin{frame}
	\frametitle{Grundkonzepte}
	\begin{itemize}
		\item Oberfläche und Eingabe mit \textit{Swing}
		\item Hauptschleife
		\item Objekte
		\begin{itemize}
			\item Bild
			\item Position
			\item Beliebige Logik
		\end{itemize}
		\item Szenen
		\begin{itemize}
			\item Mehrere Objekte
			\item Beliebige Logik
		\end{itemize}
		\item Szenenstapel
		\begin{itemize}
			\item Mehrere Szenen übereinander
		\end{itemize}
	\end{itemize}
\end{frame}

\begin{frame}
	\frametitle{Oberfläche}
	\begin{itemize}
		\item \textit{Java Swing Framework}
		\item Fenster
		\item Eingabe über \texttt{KeyListener} und \texttt{MouseListener}
	\end{itemize}
\end{frame}

\begin{frame}
	\frametitle{Hauptschleife}
	\begin{itemize}
		\item Zähler: Um seit dem letzten Durchlauf vergangene Zeit erhöhen
		\item Erreicht $\frac{1}{60}s$: Durchlauf der Logik und Grafik
	\end{itemize}
\end{frame}

\begin{frame}
	\frametitle{Objekte}
	\begin{itemize}
		\item \texttt{BaseObject}-Klasse
		\item Methoden:
		\begin{itemize}
			\item \texttt{update}
			\item \texttt{draw}
			\item \texttt{unload}
			\item \texttt{isClicked}
		\end{itemize}
	\end{itemize}
\end{frame}

\begin{frame}
	\frametitle{Szene}
	\begin{itemize}
		\item \texttt{BaseScene}-Klasse
		\item Methoden:
		\begin{itemize}
			\item \texttt{update}
			\item \texttt{draw}
			\item \texttt{unload}
		\end{itemize}
		\item Verdeckungs-Einstellungen
		\item Tags
	\end{itemize}
\end{frame}

\begin{frame}
	\frametitle{Szenenstapel}
	\begin{itemize}
		\item Liste von Szenen (\texttt{BaseScene}-Instanzen)
		\item Ausführen von \texttt{update}
		\item Ausführen von \texttt{draw}
		\item Anwenden von Verdeckungs-Einstellungen
	\end{itemize}
\end{frame}

\begin{frame}
	\frametitle{Verdeckungs-Einstellungen}
	Was soll passieren, wenn eine bestimmte Szene über einer anderen liegt?
	\vspace{0.3cm}
	
	\begin{itemize}
		\item Zuordnung Tag(s) $\rightarrow$ Regeln
		\item Werden kombiniert
		\item Soll die Szene...
		\begin{itemize}
			\item[] ... gezeichnet werden?
			\item[] ... ihre Logik ausführen?
			\item[] ... ihren Ton abspielen?
		\end{itemize}
	\end{itemize}
\end{frame}